\chapter*{Conclusion générale}
\addcontentsline{toc}{chapter}{Conclusion générale}
\markboth{Conclusion générale}{}

Ce projet de fin d'études avait pour objectif la conception et le développement d'une plateforme e-Gouvernement complète, visant à dématérialiser les services administratifs et à rapprocher les citoyens de leurs institutions. Réalisé dans le cadre d'un stage au sein de Medianet, ce travail s'est articulé autour de cinq sprints regroupés en deux releases successives, couvrant l'ensemble du cycle de vie du système, de ses fondations jusqu'à sa mise en production.

La \textbf{première release} a posé les bases du système et l'a rendu opérationnel pour les citoyens. Le \textit{Sprint 1} a mis en place le socle d'accès, la gestion des utilisateurs et des rôles ainsi que la structuration des entités administratives. Le \textit{Sprint 2} a établi le catalogue de services administratifs avec un moteur de formulaires dynamiques. Le \textit{Sprint 3} a permis la dématérialisation complète des démarches, offrant aux citoyens la possibilité de soumettre, suivre et évaluer leurs dossiers en ligne.

Cette \textbf{deuxième release} a apporté à la plateforme de nouvelles fonctionnalités liées à la communication, au suivi et à l'accompagnement des utilisateurs. Le Sprint 4, quant à lui, a permis d'introduire la gestion des réclamations, un système de notifications automatisées ainsi que des outils permettant aux administrateurs de suivre l'activité de la plateforme et de détecter d'éventuelles anomalies. Le \textit{Sprint 5} a complété l'ensemble par une vitrine institutionnelle accessible au grand public et par un assistant conversationnel intelligent fondé sur l'intelligence artificielle générative et une approche RAG, doublé d'un canal d'assistance permettant de recontacter les citoyens en difficulté.

Ce projet s'est révélé être une expérience formatrice à plusieurs égards, mêlant à la fois des défis d'ordre technique — tels que la prise en charge de formulaires dynamiques en plusieurs étapes, l'interconnexion de systèmes variés, ou encore la conception d'un assistant conversationnel RAG reposant sur une base de données vectorielle — et des problématiques métier liées à la complexité inhérente aux processus administratifs. Il m'a permis d'approfondir mes compétences en développement full-stack, en modélisation UML et en gestion de projet agile.


La plateforme pourrait néanmoins continuer à évoluer sur plusieurs points. L'ajout d'un module de paiement en ligne permettrait notamment d'étendre de façon notable l'offre de services dématérialisés. Par ailleurs, la création d'une application mobile spécifique renforcerait l'accessibilité de la solution auprès des citoyens. 

Par ailleurs, au-delà de l’assistant conversationnel déjà intégré, ses capacités pourraient être étendues afin de répondre aux besoins des administrateurs. Une assistance intelligente pourrait notamment être envisagée lors de la création des formulaires dynamiques, afin d’accompagner les utilisateurs étape par étape et de limiter les erreurs de saisie.
